\documentclass[12pt,a4paper]{report}

% =========================
% Packages
% =========================
\usepackage[a4paper,margin=2.5cm]{geometry}
\usepackage{fontspec}
\usepackage{polyglossia}
\setmainlanguage{french}

\newfontfamily\frenchfont{Times New Roman}

\usepackage{titlesec}
\usepackage{graphicx}
\usepackage{xcolor}
\usepackage{fancyhdr}
\usepackage{hyperref}
\usepackage{longtable}
\usepackage{array}
\usepackage{booktabs}
\usepackage{listings}
\usepackage{setspace}

\onehalfspacing

% =========================
% Colors
% =========================
\definecolor{maincolor}{RGB}{51,35,77}

% =========================
% Hyperref
% =========================
\hypersetup{
    colorlinks=true,
    linkcolor=maincolor,
    urlcolor=blue
}

% =========================
% Header & Footer
% =========================
\pagestyle{fancy}
\fancyhf{}
\fancyhead[R]{Résumé des logiciels open source}
\fancyhead[L]{2026}
\fancyfoot[C]{\thepage}

% =========================
% Title Style
% =========================
\titleformat{\chapter}
{\Huge\bfseries\color{maincolor}}
{\thechapter}
{1em}
{}

\titleformat{\section}
{\Large\bfseries\color{maincolor}}
{\thesection}
{1em}
{}

% =========================
% Code Style
% =========================
\lstset{
    basicstyle=\ttfamily\small,
    backgroundcolor=\color{gray!10},
    frame=single,
    breaklines=true
}

% =========================
% Document
% =========================
\begin{document}

% =========================
% Cover Page
% =========================
\begin{titlepage}
    \centering

    \vspace*{2cm}

    {\Huge\bfseries Résumé des logiciels open source \par}

    \vspace{1cm}

    

    \vspace{2cm}

    {\Large Rédigé par l’étudiante \par}

    \vspace{0.5cm}

    {\Huge Fatimetou Limam \par}

    \vfill

    {\Large Année universitaire 2026 - 2027 \par}

    \vspace{1cm}

\end{titlepage}

% =========================
% Table of Contents
% =========================
\tableofcontents
\newpage

% =========================================================
% Chapter 1
% =========================================================
\chapter{Introduction aux logiciels open source}

\section{Définition des logiciels open source}

Les logiciels open source sont des programmes dont le code source est accessible à tous pour être lu, modifié et redistribué.  
Ce modèle permet aux développeurs et aux utilisateurs de collaborer pour améliorer les logiciels.

\section{Différence entre logiciel libre et Open Source}

Il existe une différence philosophique entre le logiciel libre et l’Open Source :

\begin{itemize}
    \item Le logiciel libre met l’accent sur la liberté de l’utilisateur.
    \item L’Open Source met l’accent sur la collaboration et le développement technique.
\end{itemize}

\section{Les libertés fondamentales}

Les logiciels libres reposent sur quatre libertés fondamentales :

\begin{enumerate}
    \item La liberté d’exécuter le programme.
    \item La liberté d’étudier et de modifier le code.
    \item La liberté de redistribuer le logiciel.
    \item La liberté d’améliorer le programme et de partager les améliorations.
\end{enumerate}

\section{Histoire des logiciels open source}

La culture du partage de code a commencé dans les universités et les centres de recherche dans les années 1970.  
Elle s’est développée avec le projet GNU de Richard Stallman, puis avec Linux en 1991 créé par Linus Torvalds.

\section{Licences les plus connues}

\begin{itemize}
    \item GPL
    \item MIT
    \item Apache 2.0
    \item LGPL
\end{itemize}

\section{Avantages de l’Open Source}

\begin{itemize}
    \item Transparence
    \item Sécurité
    \item Réduction des coûts
    \item Personnalisation
    \item Développement collaboratif
\end{itemize}

\section{Modèles économiques}

Les entreprises utilisent plusieurs modèles économiques :

\begin{itemize}
    \item Open Core
    \item SaaS
    \item Services de support
    \item Double licence
\end{itemize}

\section{Défis actuels}

Parmi les principaux défis :

\begin{itemize}
    \item Cybersécurité
    \item Financement durable
    \item Gestion des communautés
    \item Concurrence des services cloud
\end{itemize}

% =========================================================
% Chapter 2
% =========================================================
\chapter{Linux et les systèmes open source}

\section{Qu’est-ce qu’un système d’exploitation ?}

Un système d’exploitation est un logiciel de base qui gère les ressources d’un ordinateur telles que la mémoire, les fichiers et les périphériques.

\section{Le système Linux}

Linux est un système d’exploitation open source créé en 1991. Il est utilisé dans les serveurs, les ordinateurs et les smartphones.

\section{Architecture de Linux}

Linux est composé de plusieurs couches :

\begin{itemize}
    \item Kernel
    \item Shell
    \item Outils GNU
    \item Interface graphique
    \item Applications
\end{itemize}

\section{Distributions Linux}

\begin{longtable}{|p{4cm}|p{4cm}|}
\hline
Distribution & Utilisation \\
\hline
Ubuntu & Adaptée aux débutants \\
\hline
Debian & Stable pour les serveurs \\
\hline
Fedora & Pour les développeurs \\
\hline
Arch Linux & Pour utilisateurs avancés \\
\hline
Kali Linux & Pour la cybersécurité \\
\hline
\end{longtable}

\section{Fichiers ISO}

Un fichier ISO est une image complète d’un système d’exploitation permettant de créer une clé USB bootable.

\section{Création d’une clé USB bootable}

Outils utilisés :

\begin{itemize}
    \item Rufus
    \item Etcher
    \item Ventoy
\end{itemize}

\section{Installation de Linux}

Étapes d’installation :

\begin{enumerate}
    \item Télécharger l’image ISO
    \item Créer une clé USB bootable
    \item Démarrer depuis la clé USB
    \item Partitionner le disque
    \item Installer le système
\end{enumerate}

\section{Gestion des paquets}

\begin{longtable}{|p{4cm}|p{5cm}|}
\hline
Gestionnaire & Distribution \\
\hline
apt & Ubuntu / Debian \\
\hline
dnf & Fedora \\
\hline
pacman & Arch Linux \\
\hline
\end{longtable}

\section{Avantages de Linux}

\begin{itemize}
    \item Gratuit
    \item Sécurisé
    \item Stable
    \item Rapide
    \item Personnalisable
\end{itemize}

% =========================================================
% Chapter 3
% =========================================================
\chapter{Guide des commandes Linux}

\section{Commandes de navigation}

\subsection*{pwd}
Affiche le répertoire courant.

\begin{lstlisting}
pwd
\end{lstlisting}

\subsection*{ls}
Affiche les fichiers et dossiers.

\begin{lstlisting}
ls -lah
\end{lstlisting}

\subsection*{cd}
Changer de répertoire.

\begin{lstlisting}
cd Documents
\end{lstlisting}

\section{Gestion des fichiers}

\subsection*{mkdir}
Créer un dossier.

\begin{lstlisting}
mkdir projet
\end{lstlisting}

\subsection*{touch}
Créer un fichier vide.

\begin{lstlisting}
touch fichier.txt
\end{lstlisting}

\subsection*{cp}
Copier des fichiers.

\begin{lstlisting}
cp fichier.txt sauvegarde.txt
\end{lstlisting}

\subsection*{mv}
Déplacer ou renommer des fichiers.

\begin{lstlisting}
mv ancien.txt nouveau.txt
\end{lstlisting}

\subsection*{rm}
Supprimer des fichiers.

\begin{lstlisting}
rm fichier.txt
\end{lstlisting}

\section{Affichage des fichiers}

\subsection*{cat}

\begin{lstlisting}
cat fichier.txt
\end{lstlisting}

\subsection*{less}

\begin{lstlisting}
less fichier.txt
\end{lstlisting}

\section{Permissions}

\subsection*{chmod}

\begin{lstlisting}
chmod 755 script.sh
\end{lstlisting}

\section{Processus}

\subsection*{ps}

\begin{lstlisting}
ps aux
\end{lstlisting}

\subsection*{top}

\begin{lstlisting}
top
\end{lstlisting}

\subsection*{kill}

\begin{lstlisting}
kill -9 PID
\end{lstlisting}

\section{Recherche}

\subsection*{find}

\begin{lstlisting}
find . -name "*.txt"
\end{lstlisting}

\subsection*{grep}

\begin{lstlisting}
grep "error" fichier.log
\end{lstlisting}

\section{Réseau}

\begin{itemize}
    \item ping
    \item curl
    \item wget
    \item ssh
    \item scp
\end{itemize}

\section{Informations système}

\begin{lstlisting}
uname -a
df -h
free -h
history
\end{lstlisting}

% =========================================================
% Conclusion
% =========================================================
% =========================================================
% Chapter 4
% =========================================================
\chapter{LibreOffice et suite bureautique open source}

LibreOffice est une suite bureautique gratuite et open source utilisée pour la création et la gestion de documents.

\section{Présentation générale}

LibreOffice permet de réaliser plusieurs tâches de bureautique comme :
\begin{itemize}
    \item la rédaction de documents,
    \item la création de tableaux,
    \item la réalisation de présentations,
    \item la gestion de données.
\end{itemize}

\section{Applications principales}

LibreOffice contient plusieurs outils :
\begin{itemize}
    \item Writer : traitement de texte,
    \item Calc : tableur,
    \item Impress : présentations,
    \item Draw : dessins et schémas,
    \item Base : bases de données,
    \item Math : formules mathématiques.
\end{itemize}

\section{Formats utilisés}

LibreOffice utilise principalement les formats OpenDocument :
\begin{itemize}
    \item .odt pour les documents,
    \item .ods pour les tableaux,
    \item .odp pour les présentations,
    \item .odg pour les dessins.
\end{itemize}
\chapter*{Conclusion}
\addcontentsline{toc}{chapter}{Conclusion}

Les logiciels open source représentent une évolution majeure dans le monde de la technologie moderne, basée sur la collaboration et la transparence.  
Linux offre un environnement puissant et stable, et ses commandes permettent un contrôle complet du système.

% =========================================================
% End
% =========================================================
\end{document}